% =====================================================================
% Mục 1. Tổng quan dữ liệu
%
% Gói cần có trong file chính: booktabs, longtable, array, caption, url (hoặc hyperref).
% Trước khi \input file này, file chính phải nạp các con số do code sinh ra:
%   % Generated by scripts/reports/report_overview.py; do not edit by hand.
\newcommand{\OvwRowsRaw}{3406}
\newcommand{\OvwColsRaw}{34}
\newcommand{\OvwDefaultNumericCols}{7}
\newcommand{\OvwDefaultObjectCols}{27}
\newcommand{\OvwDateValid}{3376}
\newcommand{\OvwDateMin}{23/03/1998}
\newcommand{\OvwDateMax}{30/09/2017}
\newcommand{\OvwDateUnknown}{30}
\newcommand{\OvwDateUnknownPct}{0,88}
\newcommand{\OvwCoreSpeedDash}{907}
\newcommand{\OvwCoreSpeedBlank}{29}
\newcommand{\OvwCoreSpeedMissing}{936}
\newcommand{\OvwCoreSpeedMissingPct}{27,48}
\newcommand{\OvwManufacturerLevels}{4}
\newcommand{\OvwNvidiaCount}{1743}
\newcommand{\OvwNvidiaPct}{51,17}
\newcommand{\OvwAmdCount}{1317}
\newcommand{\OvwAmdPct}{38,67}
\newcommand{\OvwIntelCount}{254}
\newcommand{\OvwIntelPct}{7,46}
\newcommand{\OvwAtiCount}{92}
\newcommand{\OvwAtiPct}{2,70}
\newcommand{\OvwDedicatedYes}{3018}
\newcommand{\OvwDedicatedNo}{374}
\newcommand{\OvwDedicatedMissing}{14}
\newcommand{\OvwNotebookYes}{709}
\newcommand{\OvwNotebookNo}{2697}
\newcommand{\OvwPowerObserved}{2781}
\newcommand{\OvwPowerMissing}{625}
\newcommand{\OvwPowerMissingPct}{18,35}
\newcommand{\OvwPriceMissingPct}{83,68}
\newcommand{\OvwDpMissingPct}{74,84}
\newcommand{\OvwBoostMissingPct}{57,55}

% (tạo bằng: python scripts/reports/report_overview.py; script này cũng sinh
% generated/overview_variables_rows.tex cho bảng danh sách 34 biến). Mọi con số
% trong lời văn đều là macro \Ovw..., nên chạy lại code là báo cáo tự cập nhật.
% Không tham chiếu số mục hay số bảng của mục khác (cấu trúc báo cáo chưa chốt);
% gọi theo nội dung: "phần tiền xử lý", "các mục sau".
% =====================================================================

\section{Tổng quan dữ liệu}
\label{sec:overview}

\subsection{Ngữ cảnh bài toán và nguồn dữ liệu}

Bộ xử lý đồ họa (Graphics Processing Unit -- GPU) đóng vai trò trung tâm trong xử lý đồ
họa, tính toán hiệu năng cao và trí tuệ nhân tạo. Bên cạnh hiệu năng tính toán, công suất
thiết kế tối đa (\texttt{Max\_Power}, tính bằng Watts) là thông số kỹ thuật quan trọng,
ảnh hưởng trực tiếp đến thiết kế tản nhiệt, công suất bộ nguồn khuyến nghị và chi phí vận
hành hệ thống. Trong đề tài này, nhóm thực hiện phân tích thống kê các yếu tố kiến trúc
phần cứng (quy mô đơn vị xử lý, xung nhịp, bộ nhớ, băng thông, tiến trình bán dẫn) ảnh
hưởng đến công suất thiết kế tối đa của GPU rời, đồng thời kiểm chứng ngoài mẫu trên các
thế hệ GPU giai đoạn 2018--2025.

Nguồn dữ liệu chính của phần phân tích lịch sử là tệp \texttt{All\_GPUs.csv} thuộc bộ dữ
liệu \textit{Computer Parts (CPUs and GPUs)} công bố trên Kaggle. Theo phần ghi nhận
nguồn của tác giả tập dữ liệu, thông tin được tổng hợp từ cơ sở dữ liệu phần cứng
Game-Debate, Intel và thông số kỹ thuật của các hãng sản xuất. Toàn bộ quá trình đọc,
tổng hợp và phân tích số liệu được thực hiện bằng ngôn ngữ lập trình Python.

\subsection{Mô tả tập dữ liệu gốc}

Các đặc trưng cơ bản của tệp dữ liệu gốc trước khi lọc và làm sạch gồm:
\begin{itemize}
  \item \textbf{Tên tệp tin:} \texttt{All\_GPUs.csv}
  \item \textbf{Nguồn:} {\small\url{https://www.kaggle.com/datasets/iliassekkaf/computerparts}}
  \item \textbf{Số quan trắc:} \OvwRowsRaw{} quan trắc (dòng); mỗi quan trắc là một mẫu
    sản phẩm GPU. Trong \OvwDateValid{} quan trắc có ngày phát hành hợp lệ, thời gian phát
    hành trải từ \OvwDateMin{} đến \OvwDateMax{} (còn \OvwDateUnknown{} quan trắc ghi là
    \texttt{Unknown Release Date}).
  \item \textbf{Số lượng biến:} \OvwColsRaw{} biến (cột).
  \item \textbf{Phân bố theo hãng sản xuất (\texttt{Manufacturer}):} \OvwManufacturerLevels{}
    hãng gồm Nvidia (\OvwNvidiaCount{} mẫu; \OvwNvidiaPct{}\%), AMD (\OvwAmdCount{} mẫu;
    \OvwAmdPct{}\%), Intel (\OvwIntelCount{} mẫu; \OvwIntelPct{}\%) và ATI (\OvwAtiCount{}
    mẫu; \OvwAtiPct{}\%).
  \item \textbf{Phân loại GPU:} theo biến \texttt{Dedicated}, tập dữ liệu có
    \OvwDedicatedYes{} mẫu GPU rời, \OvwDedicatedNo{} mẫu GPU tích hợp và
    \OvwDedicatedMissing{} mẫu khuyết thông tin; theo biến \texttt{Notebook\_GPU}, có
    \OvwNotebookNo{} mẫu dành cho máy để bàn và \OvwNotebookYes{} mẫu dành cho laptop.
  \item \textbf{Biến mục tiêu (\texttt{Max\_Power}):} công suất tiêu thụ tối đa (Watts),
    có \OvwPowerObserved{} quan trắc có giá trị và \OvwPowerMissing{} quan trắc khuyết
    (\OvwPowerMissingPct{}\%).
\end{itemize}

\subsection{Danh sách và phân loại các biến}

Bảng~\ref{tab:overview_variables} liệt kê toàn bộ \OvwColsRaw{} biến trong tệp gốc
\texttt{All\_GPUs.csv}, chia thành sáu nhóm theo chức năng. Trong bảng, cột ``Kiểu gốc''
ghi nhận kiểu dữ liệu khi đọc tệp CSV bằng \texttt{pandas.read\_csv} theo mặc định (khác
với bước đọc toàn bộ dưới dạng chuỗi ở phần tiền xử lý); cột ``Loại biến'' thể hiện bản
chất thống kê của từng biến; cột ``Khuyết (\%)'' được tính trên toàn bộ \OvwRowsRaw{} dòng
của tệp gốc sau khi chuẩn hóa khoảng trắng và nhận diện cả các chuỗi đại diện cho giá trị
khuyết (tỉ lệ khuyết này sẽ thay đổi ở phần tiền xử lý sau khi lọc đúng phạm vi GPU rời).

\newlength{\ovwtblw}
{\footnotesize
\setlength{\tabcolsep}{3.5pt}
\renewcommand{\arraystretch}{1.15}
\setlength{\ovwtblw}{\dimexpr\linewidth-10\tabcolsep\relax}
\begin{longtable}{@{}
    >{\ttfamily\raggedright\arraybackslash}p{0.21\ovwtblw}
    >{\centering\arraybackslash}p{0.09\ovwtblw}
    >{\raggedright\arraybackslash}p{0.20\ovwtblw}
    >{\centering\arraybackslash}p{0.11\ovwtblw}
    >{\raggedright\arraybackslash}p{0.29\ovwtblw}
    >{\raggedleft\arraybackslash}p{0.10\ovwtblw}@{}}
  \caption{Danh sách \OvwColsRaw{} biến trong tệp gốc \texttt{All\_GPUs.csv}}
  \label{tab:overview_variables} \\
  \toprule
  \normalfont\textbf{Biến} & \textbf{Kiểu gốc} & \textbf{Loại biến}
    & \textbf{Đơn vị} & \textbf{Mô tả} & \textbf{Khuyết (\%)} \\
  \midrule
  \endfirsthead

  \multicolumn{6}{@{}c@{}}{\textit{Bảng~\thetable{} -- tiếp theo}} \\
  \toprule
  \normalfont\textbf{Biến} & \textbf{Kiểu gốc} & \textbf{Loại biến}
    & \textbf{Đơn vị} & \textbf{Mô tả} & \textbf{Khuyết (\%)} \\
  \midrule
  \endhead

  \midrule
  \multicolumn{6}{@{}r@{}}{\textit{Tiếp tục ở trang sau}} \\
  \endfoot

  \bottomrule
  \multicolumn{6}{@{}p{\linewidth}@{}}{%
    \footnotesize $^{*}$ Biến không có ô khuyết mặc định (\texttt{NaN}) nhưng chứa chuỗi
    đại diện cho dữ liệu khuyết: \texttt{Release\_Date} có \OvwDateUnknown{} dòng
    ``\texttt{Unknown Release Date}''; \texttt{Core\_Speed} có \OvwCoreSpeedDash{} dòng
    ``\texttt{-}'' và \OvwCoreSpeedBlank{} dòng chỉ chứa ký tự xuống dòng.} \\
  \endlastfoot

  \csname @@input\endcsname generated/overview_variables_rows.tex
\end{longtable}
}

\subsection{Nhận xét sơ bộ về dữ liệu gốc}

{\sloppy
\begin{itemize}
  \item \textbf{Định dạng lưu trữ của các biến số:} Khi đọc bằng \texttt{pandas} theo mặc
    định, chỉ \OvwDefaultNumericCols{} trong \OvwColsRaw{} biến được nhận diện là kiểu số
    (\texttt{TMUs}, \texttt{Shader}, \texttt{Open\_GL} và bốn biến đếm cổng kết nối).
    \OvwDefaultObjectCols{} biến còn lại ở dạng chuỗi ký tự, trong đó phần lớn các biến
    định lượng (gồm cả biến mục tiêu \texttt{Max\_Power}) chứa đơn vị đo đi kèm như
    ``\texttt{141 Watts}'', ``\texttt{738 MHz}'', ``\texttt{55nm}''; một số dòng còn dùng
    ký hiệu nhân bản ``$a$ (x$N$)'' ở \texttt{ROPs} và \texttt{L2\_Cache}, hoặc
    lẫn đơn vị \texttt{GB/sec} và \texttt{MB/sec} ở \texttt{Memory\_Bandwidth}.
  \item \textbf{Dữ liệu khuyết và khuyết ẩn:} Biến mục tiêu \texttt{Max\_Power} khuyết
    \OvwPowerMissingPct{}\%. Một số biến khuyết phần lớn giá trị như
    \texttt{Release\_Price} (\OvwPriceMissingPct{}\%),
    \texttt{DisplayPort\_\allowbreak{}Connection} (\OvwDpMissingPct{}\%) và
    \texttt{Boost\_Clock} (\OvwBoostMissingPct{}\%). Ngoài ra, \texttt{Core\_Speed}
    (\OvwCoreSpeedMissingPct{}\%) và \texttt{Release\_Date} (\OvwDateUnknownPct{}\%) chứa
    giá trị khuyết dưới dạng chuỗi ký tự nên cần được nhận diện riêng thay vì chỉ kiểm tra
    \texttt{NaN}.
  \item \textbf{Phạm vi và mất cân bằng theo hãng:} Tập dữ liệu gộp cả GPU rời, GPU tích
    hợp, GPU máy để bàn, laptop và trải dài gần 20 năm (1998--2017). Theo hãng sản xuất,
    Nvidia và AMD chiếm hơn 89\% tổng số quan trắc, trong khi Intel và ATI chiếm tỉ lệ
    nhỏ. Các bước chọn biến và lọc quan sát về đúng phạm vi GPU rời được trình bày cụ thể
    ở phần tiền xử lý.
\end{itemize}
\par}
